\section{Fiber weights and spectral polynomials}\label{sec:weights}

\subsection{The class over a product of projective spaces}
Pardon's algebra $H$ is obtained from compactly supported cohomology on
spaces of cycles in families of threefolds. A closed oriented real
parameter space of dimension $b$, together with a compactly supported
$b$-cocycle on its relative cycle space, gives a class of relative
degree zero. A cocycle over an interval of such families equates its
two restrictions if its support stays compact. One may also restrict
the target to an open neighborhood containing that support.
Multiplication uses the disjoint union of targets. We use the complex
families of \cite[\S3]{Pardon} and their map $\Phi$ to the corresponding
algebra of almost-complex families \cite[\S2]{WueProduct}. The latter
allows smooth oriented parameter chains.

Take the two lines general in their Picard varieties subject to $LM=K_C$.
Thus $H^0(C,L)=H^0(C,M)=0$ and $L^2\ne K_C$.
For $A=L,M$, the vector space $H^0(C,A^i)$ has dimension
$(i-1)(g-1)$ for $i\ge2$, by Riemann--Roch and Serre duality.
A compact curve in the rank-two bundle projects finitely onto $C$,
since a fiber is affine. With multiplicities included, a cycle of
covering degree $m$ gives on each axis a monic polynomial of degree
$m$. Its coefficient of order $i$ is a holomorphic section of $A^i$.
In families these coefficients are continuous functions of the image
cycle. Near a transverse base fiber the local sheets give elementary
symmetric functions. Their extension at a branch value is the
characteristic polynomial of multiplication in the finite divisor on
the line axis. Equivalently one can push the cycle to the line and use
the coordinates on its space of monic spectral divisors
\cite[\S2]{WueBalanced}.

Fix a basis of the dual of each coefficient space using linear
combinations of evaluations at finitely many points of $C$.
Write $F\subset C$ for their union. These are $n_m$ tests on each
axis. They all vanish precisely on $m[C]$.
Let $P,Q$ denote the hyperplane lines on the factors of
$B_i=\PP^i\times\PP^{2n_m-i}$. Twist $L,M$ by $P,Q$ respectively.
Each test now takes values in $P^s$ or $Q^s$, where $s$ is its
order. Put
\[
 W_{g,m}=\prod_{s=2}^m s^{2(s-1)(g-1)}.
\]
On the relative cycle space take the pullback of the Thom class of
the direct sum of all test lines and divide by $W_{g,m}$. We denote
its image in $H$ by $Y_i$. Its cohomological degree equals
$\dim_{\RR}B_i$, and its support is the compact family of core cycles.
The choice of a general splitting computes the same normalized
one-circle classes as any splitting of the specified degrees.
Along a path in the Picard variety use a fixed sufficiently large
number of fiber evaluations and frames on small intervals. They
separate all coefficients throughout the path. Appending a redundant
test multiplies the insertion by its normalized Euler class and is
exactly the hyperplane stabilization of the Borel construction
\cite[\S2]{WueBalanced}.

\begin{lemma}\label{lem:weights}
For nonzero integral fiber weights $r,s$ the corresponding class is
\[
               X_{g,m}(r/s)=\sum_{i=0}^{2n_m}(r/s)^{i-n_m}Y_i .
\]
It is invariant under inversion of the weight ratio. The logarithmic
coefficient $P_{m,j}$ is primitive for the coproduct induced by addition
of cycles. Pullback $\rho_e$ under multiplication of cycles by $e$
sends $P_{m,j}$ to $P_{m/e,j}$ if $e\mid m$, and to zero otherwise.
\end{lemma}
\begin{proof}
Over two large projective spaces form the same supported tuple.
Pull it back by the pair of degree-$|r|,|s|$ Veronese maps from
$\PP^{2n_m}$, conjugating the corresponding source when a weight is
negative. On that oriented real-analytic base the resulting lines are
the lines of Chern classes $rH,sH$, where $H$ is the hyperplane
class. The extra scalar needed for the one-circle normalization is
$r^{-n_m}s^{-n_m}$. The fundamental class of the graph of the pair
maps to the sum of the coordinate linear products
$[\PP^i\times\PP^{2n_m-i}]$ with coefficients $r^is^{2n_m-i}$,
as is seen by pairing with the two hyperplane classes. Pairing a base
chain with the supported test cocycle commutes with its boundary:
the cocycle restrictions on common faces agree in the shifted
cochain complex for the dualizing sheaf of the base
\cite[\S2.4 and Remark~3.6]{Pardon}. These chain relations can be
taken real analytic on each simplex. Indeed embed the product of
projective spaces as Hermitian projectors, approximate a subdivided
bounding chain by affine interpolation of vertices followed by the
analytic tubular retraction, and use the straight interpolation and
retraction on the boundary prism. The support over each simplex is
the compact core family. This proves the polynomial formula. Swapping
the lines and moving the splitting along the Picard variety gives
its symmetry. For each pair of weights the cycle addition and
multiplication identities for the normalized tests are those of
\cite[\S3]{WueBalanced}; their proofs use only the orders and the
nonvanishing of the weights. Taking the logarithm and comparing
coefficients as Laurent polynomials gives the asserted operations.
\end{proof}

The degree in $z$ of the degree-$l$ logarithm is at most $n_l$ by
superadditivity of $n_l$. If $|j|>n_d$ for the largest proper divisor
$d$ of $m$, the primitive $P_{m,j}$ is therefore killed by every
proper division. By \cite[Lemma~7.1]{WueProduct} it is a finite linear
combination of $x_{h,1}$. This use of the primitive-division criterion
is independent of any evaluation by numerical invariants.

We recall the geometric detection that will determine the coefficients.
On the space of connected stable maps of a fixed arithmetic genus $k$
to fibers of a smooth family, the map to the image cycle is proper
over compactly supported insertions \cite[\S5]{WueProduct}. Pullback of
the insertion, with the oriented parameter chain, gives a class
$T_k$ in the corresponding colimit of map-space cohomology groups.
The construction annihilates products of positive-degree classes
represented on disjoint targets. In an acyclic balanced normal bundle
the only degree-one image is its zero section, so
$T_k(x_{h,1})=0$ if $k<h$. At $k=h$ the only map to this cycle is
the regular embedding, and $T_h(x_{h,1})\ne0$ by the fixed-genus
detection theorem of \cite[\S\S5--6]{WueProduct}. We shall use only
these properties of $T_k$. They require no numerical functional on a
family of singular maps.

\subsection{The algebraic pairing}
Let $\mathcal F=\overline{\mathcal M}_k(C,m)$ be the proper stack of
connected unmarked stable maps, with universal source $\pi$ and map
$f$. Write $V_A=R\pi_*f^*A$ for $A=L,M$ and
$a=k-1-m(g-1)$. Riemann--Hurwitz on the positive components gives
$a\ge0$ if $\mathcal F$ is nonempty. Both complexes have virtual
rank $-a$, and $\mathcal F$ has virtual complex dimension $2a$.
Expanding the two inverse moving Euler classes in the fiber
characters gives
\begin{equation}\label{eq:chern}
 N_{k,m,j}=
  \int_{[\mathcal F]^{\rm vir}}
    c_{a-j}(-V_L)c_{a+j}(-V_M),\qquad |j|\le a ,
\end{equation}
and $N_{k,m,j}=0$ outside this range
\cite[\S2]{BP}. The ranks refer to Euler characteristics and are
unchanged by jumps of the spaces of sections.

The pairing of the normalized tuple over $B_{n_m+j}$ with the
fiberwise algebraic stable-map virtual class is precisely
$N_{k,m,j}$. To verify it, first use the compactification
$\PP_C(\mathcal O\oplus L\oplus M)$ over a point. The tuple with
two independent torus characters $x,y$ is supported on maps into
the core and extends by zero to the projective stable-map space.
The only fixed component meeting this support is $\mathcal F$:
invariant finite multisets in the affine fiber lie at its origin.
The fixed obstruction theory there is the one for maps to $C$. The
moving theory is $V_L\otimes x\oplus V_M\otimes y$, where the
characters also denote their weight lines. Virtual localization
\cite{GP} identifies the equivariant pushforward with
\[
       x^{n_m} y^{n_m}\sum_{l}N_{k,m,l}(x/y)^l .
\]
The monomial comes from the normalized Euler class of the tests on the
fixed locus. The pushforward is an ordinary equivariant polynomial
because the support is proper. In the associated projective bundle
over $B_{n_m+j}$ the fiberwise obstruction theory is the descent of
the fixed-target theory. Keeping the residual two-torus action thus
substitutes $x+c_1(P),y+c_1(Q)$ in this polynomial. Integrating over
the base at $x=y=0$ gives the claimed coefficient, including when
$a>n_m$. This calculation uses the stable-map theory for the family
of threefolds over $B_{n_m+j}$. Maps to the total projective variety
considered as a single target have a different obstruction theory.

\subsection{Avoidance of multiples}
Assume $j>0$ by symmetry. Keep the $L$ tests equal to zero and
translate the $M$ tests to homogeneous polynomials of their respective
orders on $\PP^{n_m-j}$. Along the linear homotopy the image cycle is
confined to the $M$ axis. The prescribed coefficients determine one
relative monic divisor in each fiber, with multiplicities. Its graph
over the parameter base and the interval is compact.

Let $e>1$ divide $m$ and put $s=m/e$. The $e$th powers of monic
spectral polynomials of degree $s$ form a closed subvariety of the
affine coefficient space of codimension $n_m-n_s$. Indeed the first
$s$ coefficients of a power successively recover those of its monic
root by division by $e$. The remaining coefficients are polynomials
in the recovered coefficients, of the required orders under scaling.
This argument works also for polynomials whose divisors are reducible.
A section whose values are general homogeneous polynomials avoids the
union of the power loci if $j>n_d$: evaluation is surjective at every
point of the projective base and its dimension is less than every
such codimension. A neighborhood of the compact translated insertion
then misses all cycle multiples, since a convergent sequence of
multiples has a convergent subsequence of roots by compactness of
bounded cycles. In particular this is possible for every nonzero Laurent exponent in the stated range
when $m$ is prime. A divisor in the permitted neighborhood can still
have individual repeated components.

At the first excluded positive weight for composite $m$, take
$s=d$, $e=m/d$ and $j=n_d$. The equations on the remaining
coefficients, of orders $i>d$, are sections of $\mathcal O(i)$ with
$(i-1)(g-1)$ copies for each $i$. Their top Chern number on
$\PP^{n_m-n_d}$ is
\[
                \prod_{i=d+1}^m i^{(i-1)(g-1)} >0 .
\]
Thus the entire power locus cannot be avoided by choosing other
homogeneous values. In fact $\rho_eP_{m,n_d}=P_{d,n_d}\ne0$:
the latter has nonzero connected minimum-genus detection at its
outermost weight, as we verify next.

\begin{lemma}\label{lem:outer}
For $s\ge2$ the coefficients at the two outermost weights of
$P_s(z)=\sum_jP_{s,j}z^j$ are both
\[
 \frac{x_{H_s,1}}{A_s},\qquad H_s=1+(g-1)\binom{s+1}{2},
 \qquad A_s=\prod_{i=2}^s i^{(i-1)(g-1)} .
\]
\end{lemma}
\begin{proof}
At $j=n_s$ the second projective factor is a point.
Translate its $M$-tuple to a smooth integral spectral divisor $D$,
for example $v^s=b$ with $b$ having a simple zero and all zeros
simple. For general $M$ such a $b$ exists: the Abel map for degree
$s(g-1)\ge g$ is dominant and the reduced divisors are dense.
Eisenstein's criterion gives integrality; adjunction gives genus
$H_s$. The other tests still vanish, so the unique supported
cycle is $D$ with zero $L$ field. For the projection $f:D\to C$ the
monic equation gives $f_*\mathcal O_D=\bigoplus_{i=0}^{s-1}M^{-i}$.
Consequently $R=f^*L$ has $H^0(D,R)=0$. In the surface the normal
line $J=f^*M^s$ satisfies $JR=K_D$, so $H^1(J)=0$. Coefficient
variation gives all $H^0(J)$ of dimension $n_s$. Their tests cut
these surface deformations with degree one.
On the smooth spectral locus these are all the curves: the $M$ projection is
birational to a smooth divisor and the lift is a
section of $R$, which vanishes. This describes its topology as well,
by the family of monic divisors and the continuous inverse given by
their coefficients. The normalized $L$ tests there give the ordinary
Euler class $c_1(\mathcal O(1))^{n_s}$. Base Gysin replaces
$\PP^{2n_s}$ by a linear $\PP^{n_s}$. What remains is the point class
of the fixed smooth divisor in the coefficient space, divided by $A_s$.

We identify that point class in the universal algebra. Choose a basis
$c_1,\ldots,c_{n_s}$ of $H^1(D,R)$ and a class
$\delta\in H^1(D,RJ^{-1})$ inducing an isomorphism
$H^0(J)\to H^1(R)$ by multiplication. Indeed $H^1(RJ^{-1})$ is
dual to $H^0(J^2)$. Choose $n_s$ points whose evaluations in fixed
frames form a basis of $H^0(J)^*$, and take the sum of evaluations in
the squared frames. The pairing induced by this functional is
nondegenerate. On a cover of $C$ by two affine opens represent the
$c_i$ and $\delta$ by cocycles on their preimages in $D$. The monic
basis lifts them to polynomials in $v$ on the overlap with the
appropriate line twists. Write $P_*(v)$ for the fixed monic equation.
In the $L$ line change the transition by adding the translation
\[
  t\Big(\sum_{i=1}^{n_s}\sigma_i\widetilde c_i(v)
     +\sigma_0 P_*(v)\widetilde\delta(v)\Big),
       \qquad 0\le t\le1 ,
\]
where the $\sigma_i$ are the homogeneous coordinates on
$\PP^{n_s}$. The point test is still the equation of the fixed $M$ spectral divisor. On $D$ the translation splits for every parameter at $t=0$,
and for $t>0$ precisely at $[1:0:\cdots:0]$, with unique lift.
Thus the point test has compact support throughout this deformation.
Near the final point the lifting obstruction is a holomorphic section
of $H^1(R)\otimes\mathcal O(1)$, using reduction modulo the variable
monic polynomial to identify the cohomology in families.
Its derivative in the projective coordinates is the basis $(c_i)$
and its derivative in the divisor coordinates is multiplication by
$\delta$, up to a common sign. Both are isomorphisms. Its zero
scheme is a graph over either set of coordinates. The lifts vary
holomorphically on this graph by cohomology and base change and
$H^0(R)=0$. The divisor point test on the graph is therefore a
positive parameter point class. In that fixed target the normal
sequence of the lifted curve has subline $R$, quotient $J$, and
connecting isomorphism on $H^0(J)$ just computed. Its normal bundle
is acyclic. Deformation to the normal bundle identifies the point
class with $x_{H_s,1}$. Finally every product in the logarithmic
subtraction has strictly smaller weight than $n_s$. This proves the
formula, including its normalization.
\end{proof}
